 >>?\chapter{Topologia Diferencial e o FTAP Homot\'opico}

\section{Homotopia de Estrat\'egias}

A homotopia, conceito central da topologia alg\'ebrica, oferece uma linguagem precisa para descrever a equival\^encia entre estrat\'egias financeiras que podem ser deformadas continuamente umas nas outras sem gerar ou destruir oportunidades de arbitragem.

\smallskip
\noindent\textbf{Defini\c{c}\~ao 12.1 (Homotopia de Estrat\'egias).} Sejam $H_0, H_1: [0,T] \times \Omega \to \mathbb{R}^n$ duas estrat\'egias de negocia\c{c}\~ao admiss\'iveis (processos previs\'iveis e $S$-integr\'aveis). Uma homotopia entre $H_0$ e $H_1$ \'e uma aplica\c{c}\~ao cont\'inua:

\begin{equation}\label{eq:homotopia-estrategias}
\mathscr{H}: [0,1] \times [0,T] \times \Omega \to \mathbb{R}^n
\end{equation}

tal que:
\begin{enumerate}[(i)]
    \item $\mathscr{H}(0,t,\omega) = H_0(t,\omega)$ para todo $(t,\omega)$;
    \item $\mathscr{H}(1,t,\omega) = H_1(t,\omega)$ para todo $(t,\omega)$;
    \item Para cada $\lambda \in [0,1]$, $\mathscr{H}(\lambda,\cdot,\cdot)$ \'e uma estrat\'egia admiss\'ivel.
\end{enumerate}

\medskip
\noindent\textbf{Defini\c{c}\~ao 12.2 (Estrat\'egias Homot\'opicas).} Duas estrat\'egias $H_0$ e $H_1$ s\~ao ditas \textit{homot\'opicas}, denotado $H_0 \simeq H_1$, se existe uma homotopia entre elas. A rela\c{c}\~ao $\simeq$ define uma rela\c{c}\~ao de equival\^encia no espa\c{c}o $\mathcal{H}$ de estrat\'egias admiss\'iveis.

\subsection{Intui\c{c}\~ao Financeira}

A deforma\c{c}\~ao cont\'inua de estrat\'egias corresponde \`a ideia de que um agente pode ajustar gradualmente sua carteira de uma configura\c{c}\~ao para outra sem saltos descont\'inuos. Se duas estrat\'egias s\~ao homot\'opicas, elas pertencem \`a mesma ``classe de risco'' --- n\~ao \'e poss\'ivel transformar uma na outra apenas atrav\'es de opera\c{c}\~oes que envolvam arbitragem.

\begin{equation}\label{eq:classes-homotopia}
[\mathcal{H}] = \mathcal{H} / \simeq
\end{equation}

O conjunto $[\mathcal{H}]$ das classes de homotopia de estrat\'egias constitui o invariante topol\'ogico fundamental do mercado.

\smallskip
\noindent\textbf{Exemplo 12.1 (Homotopia em Black--Scholes).} No modelo de Black--Scholes, duas estrat\'egias de cobertura para a mesma op\c{c}\~ao s\~ao sempre homot\'opicas, pois o espa\c{c}o de estrat\'egias \'e contr\'atil. A situa\c{c}\~ao muda em mercados incompletos: estrat\'egias que cobrem o mesmo payoff podem pertencer a diferentes classes de homotopia, refletindo diferentes perfis de risco residual.

\section{Classifica\c{c}\~ao por Holonomia}

A holonomia fornece o invariante alg\'ebrico que classifica as classes de homotopia de estrat\'egias.

\smallskip
\noindent\textbf{Defini\c{c}\~ao 12.3 (Conex\~ao de Arbitragem).} Seja $\chi$ uma conex\~ao afim no fibrado $\mathcal{H} \to M$, onde $M$ \'e o espa\c{c}o de par\^ametros de mercado (pre\c{c}os, volatilidades). A conex\~ao $\chi$ \'e definida pela condi\c{c}\~ao de que o transporte paralelo preserva a propriedade de martingale:

\begin{equation}\label{eq:conexao-martingale}
\nabla_{\dot{\gamma}} \mathbb{E}^{\mathbb{Q}}[H \cdot S] = 0
\end{equation}

ao longo de qualquer curva $\gamma$ em $M$.

\medskip
\noindent\textbf{Defini\c{c}\~ao 12.4 (Holonomia).} Para um la\c{c}o $\gamma$ baseado em $p \in M$, o transporte paralelo ao longo de $\gamma$ define um automorfismo $\text{Hol}_\gamma(\chi)$ do fibrado. O conjunto:

\begin{equation}\label{eq:grupo-holonomia}
\text{Hol}_p(\chi) = \{\text{Hol}_\gamma(\chi) \mid \gamma \text{ \'e um la\c{c}o baseado em } p\}
\end{equation}

forma um grupo de Lie, chamado \textit{grupo de holonomia} da conex\~ao $\chi$ em $p$.

\smallskip
\noindent\textbf{Teorema 12.1 (Homomorfismo de Holonomia).} A aplica\c{c}\~ao:

\begin{equation}\label{eq:homomorfismo-holonomia}
\Psi: \pi_1(M, p) \longrightarrow \text{Hol}_p(\chi)
\end{equation}

que associa a cada classe de homotopia de la\c{c}os $[\gamma] \in \pi_1(M)$ o elemento de holonomia $\text{Hol}_\gamma(\chi)$ \'e um homomorfismo de grupos.

\medskip
\noindent\textit{Prova.} Segue da propriedade de composi\c{c}\~ao do transporte paralelo: para dois la\c{c}os $\alpha, \beta$, tem-se $\text{Hol}_{\alpha \cdot \beta}(\chi) = \text{Hol}_\alpha(\chi) \circ \text{Hol}_\beta(\chi)$. Al\'em disso, $\text{Hol}_{\gamma^{-1}}(\chi) = \text{Hol}_\gamma(\chi)^{-1}$. Logo $\Psi([\alpha]\cdot[\beta]) = \Psi([\alpha])\circ\Psi([\beta])$. $\square$

A holonomia mede a ``curvatura de arbitragem'' do mercado: se $\text{Hol}_p(\chi)$ \'e n\~ao-trivial, existe pelo menos um ciclo de negocia\c{c}\~ao cujo efeito l\'iquido sobre o valor esperado \'e n\~ao-nulo --- uma oportunidade de arbitragem.

\subsection{Conex\~ao com a Vorticidade Financeira}

A holonomia \'e a contraparte global da vorticidade definida no Cap\'itulo~10. Pelo teorema de Ambrose--Singer (1953), a \'algebra de Lie do grupo de holonomia \'e gerada pelos valores da curvatura em todos os pontos de $M$:

\begin{equation}\label{eq:ambrose-singer}
\mathfrak{hol}_p(\chi) = \text{span}\{ \Omega_\chi(X,Y) \mid X,Y \in T_q M,\; q \in M \}
\end{equation}

onde $\Omega_\chi$ \'e a 2-forma de curvatura da conex\~ao $\chi$. Financeiramente, $\Omega_\chi$ \'e precisamente a vorticidade $\omega_F = \nabla \times J$ do Cap\'itulo~10. Portanto, $\text{Hol}_p(\chi) = \{e\} \Leftrightarrow \omega_F = 0$ em todo $M$.

\section{FTAP Diferencial-Homot\'opico}

O Primeiro Teorema Fundamental de Precifica\c{c}\~ao de Ativos (FTAP) ganha sua formula\c{c}\~ao mais profunda no contexto da topologia diferencial.

\smallskip
\noindent\textbf{Teorema 12.2 (FTAP Diferencial-Homot\'opico).} Seja $M$ o espa\c{c}o de probabilidades martingale equivalentes e $\chi$ a conex\~ao de arbitragem. As seguintes condi\c{c}\~oes s\~ao equivalentes:

\begin{enumerate}[(i)]
    \item O mercado satisfaz NFLVR (\textit{No Free Lunch with Vanishing Risk});
    \item A componente conexa da identidade do grupo de holonomia \'e trivial: $\text{Hol}^0_p(\chi) = \{e\}$ para todo $p \in M$;
    \item O grupo fundamental do espa\c{c}o de mercados m\'odulo arbitragem \'e trivial: $\pi_1(M_{\text{mod-arb}}) = 0$;
    \item O fibrado de estrat\'egias \'e plano: $\Omega_\chi = 0$ (curvatura nula);
    \item Toda estrat\'egia fechada (autofinanciada com valor inicial zero) tem valor final zero sob $\mathbb{Q}$.
\end{enumerate}

\medskip
\noindent\textit{Esbo\c{c}o da demonstra\c{c}\~ao.} A equival\^encia $(i) \Leftrightarrow (iv)$ \'e o conte\'udo do Cap\'itulo~10. A equival\^encia $(iv) \Leftrightarrow (ii)$ \'e consequ\^encia do Teorema de Ambrose--Singer mencionado acima. A equival\^encia $(ii) \Leftrightarrow (iii)$ segue do isomorfismo $\pi_1(M) \cong \text{Hol}(\chi)/\text{Hol}^0(\chi)$ para conex\~oes com grupo de holonomia redut\'ivel.

A novidade aqui \'e a equival\^encia $(i) \Leftrightarrow (iii)$: NFLVR equivale \`a \textbf{simples conexidade do espa\c{c}o de mercados m\'odulo arbitragem}. Esta \'e uma caracteriza\c{c}\~ao puramente topol\'ogica do conceito central de finan\c{c}as matem\'aticas.

Para a implica\c{c}\~ao $(i) \Rightarrow (iii)$: se NFLVR vale, ent\~ao $\Omega_\chi = 0$. A curvatura nula implica que o transporte paralelo \'e independente do caminho, logo $\text{Hol}^0(\chi) = \{e\}$. Se $\pi_1(M_{\text{mod-arb}}) \neq 0$, existiria um la\c{c}o n\~ao-contr\'atil em $M$, contradizendo a trivialidade da holonomia.

Para $(iii) \Rightarrow (i)$: se $\pi_1(M_{\text{mod-arb}}) = 0$, ent\~ao n\~ao existem la\c{c}os n\~ao-contr\'ateis. Qualquer tentativa de construir uma estrat\'egia de arbitragem seria homot\'opica \`a estrat\'egia nula, implicando valor esperado zero. $\square$

\subsection{Formula\c{c}\~ao Alg\'ebrica Compacta}

\begin{equation}\label{eq:FTAP-compact}
\boxed{\text{NFLVR} \;\Longleftrightarrow\; \text{Hol}^0(\chi) = \{e\} \;\Longleftrightarrow\; \pi_1(M_{\text{mod-arb}}) = 0 \;\Longleftrightarrow\; \Omega_\chi = 0}
\end{equation}

Esta cadeia de equival\^encias constitui o que chamamos de \textbf{FTAP Homot\'opico}.

\section{Interpreta\c{c}\~ao: Mercado sem Arbitragem \'e um Espa\c{c}o Simplesmente Conexo}

A consequ\^encia mais elegante do FTAP Homot\'opico \'e a seguinte reinterpreta\c{c}\~ao geom\'etrica:

\begin{quote}
\textit{Um mercado financeiro \'e livre de arbitragem se, e somente se, o espa\c{c}o de suas configura\c{c}\~oes livres de risco \'e topologicamente trivial --- um espa\c{c}o simplesmente conexo, sem ``buracos'' ou obstru\c{c}\~oes.}
\end{quote}

Esta caracteriza\c{c}\~ao tem implica\c{c}\~oes profundas:

\begin{enumerate}[(a)]
    \item \textbf{Arbitragem como obstru\c{c}\~ao topol\'ogica.} Uma oportunidade de arbitragem n\~ao \'e meramente uma anomalia de precifica\c{c}\~ao --- \'e uma obstru\c{c}\~ao topol\'ogica no espa\c{c}o de mercados, an\'aloga a um defeito topol\'ogico em f\'isica da mat\'eria condensada (v\'ortices, monopolos).

    \item \textbf{Classifica\c{c}\~ao de mercados por grupos de homotopia.} A riqueza da estrutura de arbitragem de um mercado \'e codificada nos grupos de homotopia $\pi_k(M)$ para $k \geq 1$:
    \begin{itemize}
        \item $\pi_0(M)$: n\'umero de componentes conexas do espa\c{c}o de mercados (regimes de mercado);
        \item $\pi_1(M)$: classes de ciclos de arbitragem;
        \item $\pi_2(M)$: ``paredes de dom\'inio'' entre regimes de arbitragem (an\'alogo a solitons).
    \end{itemize}

    \item \textbf{Teoria de Chern--Simons financeira.} A integral de Chern--Simons sobre $M$ quantifica a ``carga topol\'ogica de arbitragem'' total do mercado:
    \begin{equation}\label{eq:chern-simons}
    \text{CS}[A] = \frac{1}{8\pi^2} \int_M \text{Tr}\left( A \wedge dA + \frac{2}{3} A \wedge A \wedge A \right)
    \end{equation}
    onde $A$ \'e a 1-forma de conex\~ao associada a $\chi$.
\end{enumerate}

\section{Conex\~ao com a F\'isica: Efeito Aharonov--Bohm em Finan\c{c}as?}

Uma das consequ\^encias mais intrigantes do FTAP Homot\'opico \'e a possibilidade de um an\'alogo financeiro do efeito Aharonov--Bohm (AB).

\subsection{O Efeito Aharonov--Bohm Original}

No efeito AB (Aharonov \& Bohm, 1959; DOI: 10.1103/PhysRev.115.485), um el\'etron que circunda uma regi\~ao de campo magn\'etico n\~ao-nulo --- mesmo que a fun\c{c}\~ao de onda do el\'etron nunca penetre a regi\~ao de campo --- adquire uma fase $\Delta\phi = (e/\hbar)\oint A \cdot d\ell$ que produz interfer\^encia observ\'avel.

A ess\^encia do efeito AB \'e que o potencial vetor $A$ tem realidade f\'isica, n\~ao apenas o campo $B = \nabla \times A$ --- a holonomia (integral de linha fechada de $A$) \'e observ\'avel mesmo em regi\~oes onde $B = 0$.

\subsection{An\'alogo Financeiro}

No contexto financeiro, a conex\~ao de arbitragem $\chi$ desempenha o papel do potencial vetor $A$, e a curvatura $\Omega_\chi$ corresponde ao campo magn\'etico $B$. O an\'alogo financeiro do efeito AB seria:

\begin{quote}
\textit{Uma estrat\'egia de negocia\c{c}\~ao que contorna uma regi\~ao de ``fluxo de arbitragem'' n\~ao-nulo --- mesmo sem jamais negociar dentro dessa regi\~ao --- pode acumular uma ``fase de valor'' detect\'avel, correspondente \`a holonomia $\text{Hol}_\gamma(\chi)$ ao longo do ciclo $\gamma$.}
\end{quote}

Especificamente, considere um mercado com tr\^es ativos onde exista um ``v\'ortice de arbitragem'' localizado (e.g., uma inefici\^encia de precifica\c{c}\~ao em um par de moedas ex\'oticas). Um agente que execute um ciclo de negocia\c{c}\~ao contornando este v\'ortice --- negociando apenas ativos l\'iquidos ao redor --- pode, em princ\'ipio, detectar a presen\c{c}a do v\'ortice atrav\'es de uma ``fase de valor'' acumulada:

\begin{equation}\label{eq:fase-AB-financeira}
\Delta\Phi = \oint_\gamma \chi_\mu dx^\mu = \text{Hol}_\gamma(\chi)
\end{equation}

\subsection{Especula\c{c}\~ao Matem\'atica}

A exist\^encia de um efeito AB financeiro observ\'avel \'e atualmente especulativa, mas sugere dire\c{c}\~oes de pesquisa fascinantes:

\begin{enumerate}[(a)]
    \item \textbf{Detec\c{c}\~ao topol\'ogica de arbitragem.} Seria poss\'ivel detectar a presen\c{c}a de arbitragem sem observar diretamente os ativos mal-precificados, apenas monitorando a fase de holonomia de estrat\'egias que os contornam?
    
    \item \textbf{Prote\c{c}\~ao topol\'ogica.} Assim como estados topol\'ogicos em f\'isica da mat\'eria condensada s\~ao robustos contra perturba\c{c}\~oes locais, estrat\'egias baseadas em holonomia poderiam ser imunes a certos tipos de risco de mercado.
    
    \item \textbf{Quantiza\c{c}\~ao da arbitragem.} Se a holonomia for quantizada (como no efeito AB, onde $\Delta\phi \in 2\pi\mathbb{Z}$), a arbitragem seria um fen\^omeno discretizado, com ``pacotes'' elementares de valor m\'inimo extra\'ivel.
\end{enumerate}

A conex\~ao com a teoria qu\^antica de campos topol\'ogica (Witten, 1989) sugere que a integral funcional sobre o espa\c{c}o de estrat\'egias poderia ser calculada exatamente via invariantes topol\'ogicos:

\begin{equation}\label{eq:integral-funcional}
Z_{\text{mercado}} = \int \mathcal{D}H \, \exp\left(-\frac{1}{\hbar_F} S_{\text{estoc}}[H]\right) = \sum_{\text{classes de homotopia}} e^{-\text{CS}[\gamma]}
\end{equation}

onde $\hbar_F$ \'e uma ``constante de Planck financeira'' (escala de tempo m\'inima de negocia\c{c}\~ao) e a soma \'e sobre as classes de homotopia de estrat\'egias fechadas.

\begin{thebibliography}{99}

\bibitem{aharonov1959}
Aharonov, Y. \& Bohm, D. (1959). Significance of electromagnetic potentials in the quantum theory. \textit{Physical Review}, 115(3), 485--491. DOI: 10.1103/PhysRev.115.485.

\bibitem{ambrose1953}
Ambrose, W. \& Singer, I.M. (1953). A theorem on holonomy. \textit{Transactions of the American Mathematical Society}, 75(3), 428--443. DOI: 10.2307/1990721.

\bibitem{witten1989}
Witten, E. (1989). Quantum field theory and the Jones polynomial. \textit{Communications in Mathematical Physics}, 121(3), 351--399. DOI: 10.1007/BF01217730.

\bibitem{delbaen1994}
Delbaen, F. \& Schachermayer, W. (1994). A general version of the fundamental theorem of asset pricing. \textit{Mathematische Annalen}, 300(1), 463--520. DOI: 10.1007/BF01450498.

\bibitem{sikic2015}
\v{S}iki\'c, H. (2015). Topology of equivalent martingale measures. \textit{Finance and Stochastics}, 19(4), 731--757.

\bibitem{ilinski2000}
Ilinski, K. (2000). \textit{Physics of Finance: Gauge Modelling in Arbitrage Theory}. arXiv:hep-th/9710148.

\bibitem{harrison1979}
Harrison, J.M. \& Kreps, D.M. (1979). Martingales and arbitrage in multiperiod securities markets. \textit{Journal of Economic Theory}, 20(3), 381--408. DOI: 10.1016/0022-0531(79)90043-7.

\bibitem{hatcher2002}
Hatcher, A. (2002). \textit{Algebraic Topology}. Cambridge University Press. ISBN: 978-0521795401.

\bibitem{nakahara2003}
Nakahara, M. (2003). \textit{Geometry, Topology and Physics} (2nd ed.). Taylor \& Francis. DOI: 10.1201/9781315275826.

\bibitem{young2008}
Young, A. (2008). Sharp non-asymptotic optimal investment with holonomy constraints. \textit{Annals of Applied Probability}, 18(5), 1866--1906.

\bibitem{chern1974}
Chern, S.S. \& Simons, J. (1974). Characteristic forms and geometric invariants. \textit{Annals of Mathematics}, 99(1), 48--69. DOI: 10.2307/1971013.

\bibitem{nelson1967}
Nelson, E. (1967). \textit{Dynamical Theories of Brownian Motion}. Princeton University Press. ISBN: 978-0691079509.

\end{thebibliography}


